\documentclass[11pt,letterpaper]{article}
\def\SuiteDocTitle{Prerequisite-Based Learning Curriculum}
\def\SuiteDocCode{CU-01}
% Shared professional style: pdfLaTeX compatible, no external font files.
\usepackage[T1]{fontenc}
\usepackage[utf8]{inputenc}
\usepackage{lmodern}
\usepackage{microtype}
\usepackage[margin=0.83in,headheight=16pt,footskip=28pt]{geometry}
\usepackage[table]{xcolor}
\definecolor{navy}{HTML}{12304A}
\definecolor{teal}{HTML}{087F83}
\definecolor{gold}{HTML}{B78633}
\definecolor{blueback}{HTML}{F0F5F8}
\definecolor{tealback}{HTML}{EEF8F6}
\definecolor{amberback}{HTML}{FFF7E8}
\definecolor{grayback}{HTML}{F6F7F8}
\usepackage{array,tabularx,longtable,booktabs}
\usepackage{enumitem}
\usepackage{amssymb}
\usepackage{titlesec}
\usepackage{fancyhdr,lastpage}
\usepackage{tcolorbox}
\tcbuselibrary{breakable}
\usepackage{pgffor}
\usepackage{hyperref}
\hypersetup{colorlinks=true,linkcolor=teal,urlcolor=teal,pdfauthor={Engineering Learning Series},pdftitle={\SuiteDocTitle}}
\setlength{\parindent}{0pt}
\setlength{\parskip}{5.5pt}
\setlist{topsep=3pt}
\setcounter{tocdepth}{2}
\titleformat{\section}{\Large\bfseries\color{navy}}{\thesection}{0.7em}{}
\titleformat{\subsection}{\large\bfseries\color{teal}}{\thesubsection}{0.6em}{}
\titleformat{\subsubsection}{\normalsize\bfseries\color{navy}}{\thesubsubsection}{0.6em}{}
\pagestyle{fancy}
\fancyhf{}
\fancyhead[L]{\scriptsize\color{navy}AGILE REQUIREMENTS ENGINEERING}
\fancyhead[R]{\scriptsize\color{teal}\SuiteDocCode}
\fancyfoot[L]{\scriptsize Independent learning and implementation materials}
\fancyfoot[R]{\scriptsize Page \thepage\ of \pageref{LastPage}}
\renewcommand{\headrulewidth}{0.35pt}
\renewcommand{\footrulewidth}{0pt}
\newcommand{\SuiteHeader}[2]{%
{\color{teal}\small\bfseries ENGINEERING DOCUMENTATION SUITE \hfill #2}\par\vspace{8pt}
{\color{navy}\LARGE\bfseries\SuiteDocTitle}\par\vspace{5pt}
{\color{black!68}\normalsize #1}\par\vspace{8pt}
{\color{gold}\rule{\linewidth}{2pt}}\par\vspace{8pt}}
\newcommand{\WritingLines}[1]{\par\foreach \n in {1,...,#1}{\noindent\rule{\linewidth}{0.35pt}\par\vspace{0.28cm}}}
\newcommand{\FreeSpace}[1]{\begin{tcolorbox}[colback=white,colframe=black!25,height=#1,boxrule=0.5pt]\end{tcolorbox}}
\newcommand{\Field}[1]{\noindent\textbf{#1}\quad\rule{0.42\linewidth}{0.35pt}\par\vspace{0.3cm}}
\newcommand{\CheckItem}[1]{\noindent$\square$\quad #1\par\vspace{2pt}}
\newcommand{\SmallNote}[1]{{\footnotesize\color{black!62}#1}\par}
\newcommand{\Guidance}[1]{\begin{tcolorbox}[colback=blueback,colframe=teal,breakable,title=Guidance]#1\end{tcolorbox}}
\sloppy

\begin{document}
\SuiteHeader{Twelve-module project-driven path | competency evidence and capstone}{Reference}
\tableofcontents\medskip\noindent\rule{\linewidth}{0.4pt}\medskip
\section{Program rationale}
This curriculum is prerequisite-based, project-driven and organized into 12 modules. It follows the book’s emphasis on conversational discovery, user-valued story units, confirmation by tests, relative estimation and replanning. The path is not merely a sequential list of chapters: case-study chapters are revisited as applied reinforcement.

\section{Learning objectives}
\begin{itemize}[leftmargin=*,itemsep=2pt]
\item Elicit goals from real roles and distinguish evidence from interpretation.
\item Write, refine, size and prioritize demonstrable user-valued increments.
\item Derive testable acceptance rules and cross-cutting constraints.
\item Create defensible release and iteration forecasts without false precision.
\item Maintain enough durable traceability for a secure, maintainable delivery process.
\end{itemize}
\section{Prerequisite map}
\begin{longtable}{|p{0.12\linewidth}|p{0.18\linewidth}|p{0.12\linewidth}|p{0.18\linewidth}|p{0.25\linewidth}|}\hline
\textbf{Module} & \textbf{Focus} & \textbf{Depends on} & \textbf{Read chapters} & \textbf{Build / assessment} \\ \hline\endfirsthead\hline
\textbf{Module} & \textbf{Focus} & \textbf{Depends on} & \textbf{Read chapters} & \textbf{Build / assessment} \\ \hline\endhead
1 & Requirements as a communication discipline & None & 1, 12, 13 & Role-aware story; artifact choice note \\ \hline
2 & Role discovery and empirical research & 1 & 3, 4, 5, 17 & Role cards and interview plan \\ \hline
3 & Story writing and quality & 1–2 & 2, 7, 14, 18 & Six revised stories and scored examples \\ \hline
4 & Vertical decomposition and readiness & 3 & 7, 10, 18 & Epic slicing map \\ \hline
5 & Acceptance engineering & 3–4 & 6, 21 & Eight acceptance cases \\ \hline
6 & Relative estimation & 3–5 & 8, 19 & Estimated sample backlog \\ \hline
7 & Release planning and risk & 6 & 9, 20 & Release scenarios \\ \hline
8 & Iteration execution and metrics & 6–7 & 10, 11 & Iteration goal and evidence ledger \\ \hline
9 & Scrum integration and defects & 5–8 & 15, 16 & Refinement operating procedure \\ \hline
10 & Durable requirements and quality constraints & 5, 9 & 12, 13, 16 & Change-impact and constraints \\ \hline
11 & Enterprise governance and review & 7–10 & 5, 11, 14, 16 & Traceability and governance review \\ \hline
12 & Integrated capstone & All prior modules & 17–21 (revisit) & Full capstone evidence bundle \\ \hline
\end{longtable}
\section{Teaching sequence}
\subsection{Module 1: Requirements as a communication discipline}
Prerequisite modules: None. Primary reading: book Chapters 1, 12, 13. Learning evidence: Role-aware story; artifact choice note. Recommended laboratory IDs: 02.

\begin{itemize}[leftmargin=*,itemsep=2pt]
\item Begin with a short explanation and an example from the running service-request case.
\item Complete one collaborative exercise that produces a durable artifact.
\item Peer-review assumptions, observable outcomes and security-sensitive gaps.
\item Record a decision, correction or explicit question before moving forward.
\end{itemize}
\subsection{Module 2: Role discovery and empirical research}
Prerequisite modules: 1. Primary reading: book Chapters 3, 4, 5, 17. Learning evidence: Role cards and interview plan. Recommended laboratory IDs: 01, 04.

\begin{itemize}[leftmargin=*,itemsep=2pt]
\item Begin with a short explanation and an example from the running service-request case.
\item Complete one collaborative exercise that produces a durable artifact.
\item Peer-review assumptions, observable outcomes and security-sensitive gaps.
\item Record a decision, correction or explicit question before moving forward.
\end{itemize}
\subsection{Module 3: Story writing and quality}
Prerequisite modules: 1–2. Primary reading: book Chapters 2, 7, 14, 18. Learning evidence: Six revised stories and scored examples. Recommended laboratory IDs: 02.

\begin{itemize}[leftmargin=*,itemsep=2pt]
\item Begin with a short explanation and an example from the running service-request case.
\item Complete one collaborative exercise that produces a durable artifact.
\item Peer-review assumptions, observable outcomes and security-sensitive gaps.
\item Record a decision, correction or explicit question before moving forward.
\end{itemize}
\subsection{Module 4: Vertical decomposition and readiness}
Prerequisite modules: 3. Primary reading: book Chapters 7, 10, 18. Learning evidence: Epic slicing map. Recommended laboratory IDs: 03.

\begin{itemize}[leftmargin=*,itemsep=2pt]
\item Begin with a short explanation and an example from the running service-request case.
\item Complete one collaborative exercise that produces a durable artifact.
\item Peer-review assumptions, observable outcomes and security-sensitive gaps.
\item Record a decision, correction or explicit question before moving forward.
\end{itemize}
\subsection{Module 5: Acceptance engineering}
Prerequisite modules: 3–4. Primary reading: book Chapters 6, 21. Learning evidence: Eight acceptance cases. Recommended laboratory IDs: 05.

\begin{itemize}[leftmargin=*,itemsep=2pt]
\item Begin with a short explanation and an example from the running service-request case.
\item Complete one collaborative exercise that produces a durable artifact.
\item Peer-review assumptions, observable outcomes and security-sensitive gaps.
\item Record a decision, correction or explicit question before moving forward.
\end{itemize}
\subsection{Module 6: Relative estimation}
Prerequisite modules: 3–5. Primary reading: book Chapters 8, 19. Learning evidence: Estimated sample backlog. Recommended laboratory IDs: 06.

\begin{itemize}[leftmargin=*,itemsep=2pt]
\item Begin with a short explanation and an example from the running service-request case.
\item Complete one collaborative exercise that produces a durable artifact.
\item Peer-review assumptions, observable outcomes and security-sensitive gaps.
\item Record a decision, correction or explicit question before moving forward.
\end{itemize}
\subsection{Module 7: Release planning and risk}
Prerequisite modules: 6. Primary reading: book Chapters 9, 20. Learning evidence: Release scenarios. Recommended laboratory IDs: 07.

\begin{itemize}[leftmargin=*,itemsep=2pt]
\item Begin with a short explanation and an example from the running service-request case.
\item Complete one collaborative exercise that produces a durable artifact.
\item Peer-review assumptions, observable outcomes and security-sensitive gaps.
\item Record a decision, correction or explicit question before moving forward.
\end{itemize}
\subsection{Module 8: Iteration execution and metrics}
Prerequisite modules: 6–7. Primary reading: book Chapters 10, 11. Learning evidence: Iteration goal and evidence ledger. Recommended laboratory IDs: 06.

\begin{itemize}[leftmargin=*,itemsep=2pt]
\item Begin with a short explanation and an example from the running service-request case.
\item Complete one collaborative exercise that produces a durable artifact.
\item Peer-review assumptions, observable outcomes and security-sensitive gaps.
\item Record a decision, correction or explicit question before moving forward.
\end{itemize}
\subsection{Module 9: Scrum integration and defects}
Prerequisite modules: 5–8. Primary reading: book Chapters 15, 16. Learning evidence: Refinement operating procedure. Recommended laboratory IDs: 08.

\begin{itemize}[leftmargin=*,itemsep=2pt]
\item Begin with a short explanation and an example from the running service-request case.
\item Complete one collaborative exercise that produces a durable artifact.
\item Peer-review assumptions, observable outcomes and security-sensitive gaps.
\item Record a decision, correction or explicit question before moving forward.
\end{itemize}
\subsection{Module 10: Durable requirements and quality constraints}
Prerequisite modules: 5, 9. Primary reading: book Chapters 12, 13, 16. Learning evidence: Change-impact and constraints. Recommended laboratory IDs: 08.

\begin{itemize}[leftmargin=*,itemsep=2pt]
\item Begin with a short explanation and an example from the running service-request case.
\item Complete one collaborative exercise that produces a durable artifact.
\item Peer-review assumptions, observable outcomes and security-sensitive gaps.
\item Record a decision, correction or explicit question before moving forward.
\end{itemize}
\subsection{Module 11: Enterprise governance and review}
Prerequisite modules: 7–10. Primary reading: book Chapters 5, 11, 14, 16. Learning evidence: Traceability and governance review. Recommended laboratory IDs: 08.

\begin{itemize}[leftmargin=*,itemsep=2pt]
\item Begin with a short explanation and an example from the running service-request case.
\item Complete one collaborative exercise that produces a durable artifact.
\item Peer-review assumptions, observable outcomes and security-sensitive gaps.
\item Record a decision, correction or explicit question before moving forward.
\end{itemize}
\subsection{Module 12: Integrated capstone}
Prerequisite modules: All prior modules. Primary reading: book Chapters 17–21 (revisit). Learning evidence: Full capstone evidence bundle. Recommended laboratory IDs: 09.

\begin{itemize}[leftmargin=*,itemsep=2pt]
\item Begin with a short explanation and an example from the running service-request case.
\item Complete one collaborative exercise that produces a durable artifact.
\item Peer-review assumptions, observable outcomes and security-sensitive gaps.
\item Record a decision, correction or explicit question before moving forward.
\end{itemize}
\section{Mastery thresholds and portfolio}
\begin{longtable}{|p{0.28\linewidth}|p{0.60\linewidth}|}\hline
\textbf{Level} & \textbf{Evidence} \\ \hline\endfirsthead\hline
\textbf{Level} & \textbf{Evidence} \\ \hline\endhead
Foundation & Role model, user stories and clear explanation of Card–Conversation–Confirmation \\ \hline
Practitioner & Vertically sliced backlog, acceptance examples and a defended estimate \\ \hline
Advanced & Risk-aware release plan, traceability, tests and change impact \\ \hline
Portfolio-ready & Capstone can be followed and challenged by an independent reviewer \\ \hline
\end{longtable}
Assess with a mix of peer demonstration, scenario review, completed labs and scorecards. Do not use arbitrary story point totals as proof of engineering proficiency.

\section{Study cadence}
A suggested weekly rhythm is two reading sessions, a workshop, a laboratory and a short retrospective. At minimum, capture what changed in your understanding, which assumption was invalidated, and what evidence you will carry into the next module.


\end{document}
